\documentclass[10pt,a4paper]{amsart}
\usepackage{amsmath,amssymb,amsthm,graphicx}
\usepackage{verbatim}
\newenvironment{reference}{\comment}{\endcomment}
%\newenvironment{reference}{}{}
\newenvironment{slogan}{\comment}{\endcomment}
\newenvironment{history}{\comment}{\endcomment}

% For commutative diagrams we use Xy-pic
\usepackage[all]{xy}

% We use 2cell for 2-commutative diagrams.
\xyoption{2cell}
% Theorem environments.
%
\theoremstyle{plain}
\newtheorem{theorem}[subsection]{Theorem}
\newtheorem{proposition}[subsection]{Proposition}
\newtheorem{lemma}[subsection]{Lemma}

\theoremstyle{definition}
\newtheorem{definition}[subsection]{Definition}
\newtheorem{example}[subsection]{Example}
\newtheorem{exercise}[subsection]{Exercise}
\newtheorem{situation}[subsection]{Situation}

\theoremstyle{remark}
\newtheorem{remark}[subsection]{Remark}
\newtheorem{remarks}[subsection]{Remarks}

\numberwithin{equation}{subsection}

% Macros
%
\def\lim{\mathop{\mathrm{lim}}\nolimits}
\def\colim{\mathop{\mathrm{colim}}\nolimits}
\def\Spec{\mathop{\mathrm{Spec}}}
\def\Hom{\mathop{\mathrm{Hom}}\nolimits}
\def\Ext{\mathop{\mathrm{Ext}}\nolimits}
\def\SheafHom{\mathop{\mathcal{H}\!\mathit{om}}\nolimits}
\def\SheafExt{\mathop{\mathcal{E}\!\mathit{xt}}\nolimits}
\def\Sch{\mathit{Sch}}
\def\Mor{\mathop{\mathrm{Mor}}\nolimits}
\def\Ob{\mathop{\mathrm{Ob}}\nolimits}
\def\Sh{\mathop{\mathit{Sh}}\nolimits}
\def\NL{\mathop{N\!L}\nolimits}
\def\CH{\mathop{\mathrm{CH}}\nolimits}
\def\proetale{{pro\text{-}\acute{e}tale}}
\def\etale{{\acute{e}tale}}
\def\QCoh{\mathit{QCoh}}
\def\Ker{\mathop{\mathrm{Ker}}}
\def\Im{\mathop{\mathrm{Im}}}
\def\Coker{\mathop{\mathrm{Coker}}}
\def\Coim{\mathop{\mathrm{Coim}}}

% Boxtimes
%
\DeclareMathSymbol{\boxtimes}{\mathbin}{AMSa}{"02}

%
% Macros for moduli stacks/spaces
%
\def\QCohstack{\mathcal{QC}\!\mathit{oh}}
\def\Cohstack{\mathcal{C}\!\mathit{oh}}
\def\Spacesstack{\mathcal{S}\!\mathit{paces}}
\def\Quotfunctor{\mathrm{Quot}}
\def\Hilbfunctor{\mathrm{Hilb}}
\def\Curvesstack{\mathcal{C}\!\mathit{urves}}
\def\Polarizedstack{\mathcal{P}\!\mathit{olarized}}
\def\Complexesstack{\mathcal{C}\!\mathit{omplexes}}
% \Pic is the operator that assigns to X its picard group, usage \Pic(X)
% \Picardstack_{X/B} denotes the Picard stack of X over B
% \Picardfunctor_{X/B} denotes the Picard functor of X over B
\def\Pic{\mathop{\mathrm{Pic}}\nolimits}
\def\Picardstack{\mathcal{P}\!\mathit{ic}}
\def\Picardfunctor{\mathrm{Pic}}
\def\Deformationcategory{\mathcal{D}\!\mathit{ef}}
\usepackage{cleveref}
\crefname{originallemma}{original lemma}{original lemmas}
\Crefname{originallemma}{Original lemma}{Original lemmas}
\makeatletter
\newlabel{more-morphisms-lemma-flat-locus-base-change}{{047C}{source}{}{sourcepointer-1}{}}
\newlabel{more-morphisms-lemma-flat-locus-base-change@cref}{{[originallemma][0][]047C}{[1][0][]source}{}{}{}}
\makeatother
\title{Flatness of a module at a point under base change}
\author{The Stacks Project Authors}
\date{stacks-project:spaces-morphisms.tex:local-snapshot-20261012; GNU FDL 1.2 or later}
\begin{document}
\maketitle
\setcounter{subsection}{1}
\noindent\textit{Editorial note.} All delivered mathematical text is the author's own, copied from the pinned Stacks Project chapter file spaces-morphisms.tex; no statement or proof step was rewritten, reordered or removed.\par\medskip
\noindent\textit{Editorial note.} The statement displays the cartesian square and includes the 'in particular' sentence; the proof handles both numbered assertions at once through the same presentation.\par\medskip
\noindent\textit{Editorial note.} The single reference of the counted spans is recorded as a non-load-bearing navigation pointer; the argument that consumes it stays inside the counted proof span.\par\medskip
\noindent\textit{Editorial note.} Preamble: the three native excerpt spans of sources/stacks-spaces-morphisms/preamble.tex (lines 8-13, 14-18 and 154-217). The chapter file itself carries no preamble text.\par\medskip
\noindent\textit{Allowed standard prerequisites (editorial).}\begin{itemize}
\item Testing flatness of a quasi-coherent module at a point through an etale presentation. used at proof lines 6444-6447. Replaces the algebraic space statement by the scheme statement for U' over V'.
\item Base change behaviour of the flat locus for morphisms of schemes. More on Morphisms, Lemma more-morphisms-lemma-flat-locus-base-change, used at proof lines 6447-6448. Is the scheme input that settles both assertions.
\item Existence of affine etale presentations over a fixed algebraic space. used at proof lines 6440-6444. Produces V, U, V' and the surjective etale morphism U' -> X'.
\end{itemize}\medskip
\ifdefined\hypertarget\hypertarget{sourcepointer-1}{}\fi
\noindent\textit{Source reference.} Original lemma 047C: focused-seeds/AGP/sources/stacks-snapshot-20261011/tags.txt, source lines 5503--5503. Navigation pointer into the Stacks Project; the standard fact it supplies is applied inside the counted proof span and the argument resting on it stays there.\par\medskip
\begin{lemma}
\label{lemma-base-change-module-flat}
Let $S$ be a scheme. Let
$$
\xymatrix{
X' \ar[d]_{f'} \ar[r]_{g'} & X \ar[d]^f \\
Y' \ar[r]^g & Y
}
$$
be a cartesian diagram of algebraic spaces over $S$. Let $x' \in |X'|$
with image $x \in |X|$. Let $\mathcal{F}$ be a quasi-coherent
sheaf on $X$ and denote $\mathcal{F}' = (g')^*\mathcal{F}$.
\begin{enumerate}
\item If $\mathcal{F}$ is flat at $x$ over $Y$
then $\mathcal{F}'$ is flat at $x'$ over $Y'$.
\item If $g$ is flat at $f'(x')$ and
$\mathcal{F}'$ is flat at $x'$ over $Y'$, then
$\mathcal{F}$ is flat at $x$ over $Y$.
\end{enumerate}
In particular, if $\mathcal{F}$ is flat over $Y$, then
$\mathcal{F}'$ is flat over $Y'$.
\end{lemma}

\begin{proof}
Choose a scheme $V$ and a surjective \'etale morphism $V \to Y$.
Choose a scheme $U$ and a surjective \'etale morphism $U \to V \times_Y X$.
Choose a scheme $V'$ and a surjective \'etale morphism $V' \to V \times_Y Y'$.
Then $U' = V' \times_V U$ is a scheme endowed with a surjective \'etale
morphism $U' = V' \times_V U \to Y' \times_Y X = X'$. Pick $u' \in U'$
mapping to $x' \in |X'|$. Then we can check flatness of
$\mathcal{F}'$ at $x'$ over $Y'$ in terms of flatness of
$\mathcal{F}'|_{U'}$ at $u'$ over $V'$. Hence the lemma follows from
More on Morphisms, Lemma \ref{more-morphisms-lemma-flat-locus-base-change}.
\end{proof}

\end{document}
